% GENERATED_BY: oc_core_monograph_translation_draft_pipeline_v1.py
% AI_ASSISTED_TRANSLATION_DRAFT: TRUE
% THEOREM_REVIEW_STATUS: PENDING
% THEOREM_REVIEW_MODE: FORMAL_AI_GATE__CODEX_CLI_CHATGPT
% THEOREM_REVIEW_REVIEWER_ID:
% THEOREM_REVIEW_AUTHORITY_CLASS:
% THEOREM_REVIEWED_AT:
% THEOREM_REVIEW_DOSSIER_REF:
% SOURCE_LANGUAGE: EN
% TARGET_LANGUAGE: RU
% SOURCE_EN_REF: content/k_levels/klevels_master.tex
% ================================================================
% ==== FILE: content/k_levels/klevels_master.tex
% ================================================================



\paragraph{Canonical tuple projection note.}
Local K-level displays in this module are projections of the canonical public tuple \(K=(\Omega,\partial\Omega,A,\Theta,P,J,C,k,M)\). When a display omits \(M\), reorders components, or uses level-indexed shorthand, the omitted embedding context is inherited from the surrounding level discussion rather than introduced as a competing definition.

\section{Расширенные модули и межуровневые структуры}
\label{sec:klevels-master}

Онтология континуумов организована как иерархическая последовательность
уровней
\[
K_0, K_1, \dots, K_{12},
\]
каждый из которых представляет качественно отличную форму структуры,
динамики и допустимых состояний.
Каждый уровень $K_x$ имеет:
\begin{itemize}
    \item пространство состояний $\Omega(K_x)$,
    \item множество осей $A(K_x)$,
    \item потенциалы $P(K_x)$,
    \item пороги $\Theta(K_x)$,
    \item потоки $J(K_x)$,
    \item циклы $C(K_x)$ с характерными временами $\tau(K_x)$,
    \item меру континуума $k(K_x)$,
    \item и границу $\partial\Omega(K_x)$, описывающую условия
          коллапса или перехода.
\end{itemize}

Каждый $K_x$ вложен в соответствующее мета-пространство $M_x$, которое
задаёт допустимые оси и состояния:
\[
A(K_x) \subseteq A(M_x), \qquad
\Omega(K_x) \subseteq \Omega(M_x).
\]
Мета-пространство ограничивает число степеней свободы континуума и
тем самым определяет его эволюцию.


% ================================================================
\subsection{Размерностный рост и переходы $K_x \to K_{x+1}$}
% ================================================================

Переход от уровня $K_x$ к $K_{x+1}$ происходит, когда:
\begin{enumerate}
    \item появляется новый класс различий, который нельзя представить
          в пределах существующих осей $A(K_x)$;
    \item структурное напряжение превышает порог:
    \[
        T(K_x) > \Theta_{\mathrm{dim}}(K_x);
    \]
    \item мета-пространство требует более высокоразмерного представления;
    \item существующее пространство конфигураций $\Omega(K_x)$ становится
          недостаточным или схлопывается на свою границу $\partial\Omega(K_x)$.
\end{enumerate}

Процесс необратим:
\[
\dim K_{x+1} \ge \dim K_x,
\]
а любая «редукция» размерности соответствует смерти континуума:
\[
\Omega(K_x) = \varnothing.
\]


% ================================================================
\subsection{Единая формальная структура уровней}
% ================================================================

Несмотря на качественное разнообразие уровней, каждый $K_x$ обладает общей
формальной структурой.
Эволюция континуума задаётся уравнением:
\[
K_x(t + dt) = E\bigl(K_x(t), M_x(t)\bigr),
\]
где $E$ — оператор эволюции, ограниченный мета-пространством
$M_x$.
Континуум эволюционирует только внутри допустимой области своего
мета-пространства:
\[
K_x(t+dt) \in \Omega(M_x)
\quad\text{или эволюция запрещена}.
\]

Мера континуума задаётся как:
\[
k(K_x) =
k_{\Omega} \cdot k_{A} \cdot k_{C} \cdot k_{J} \cdot k_{T},
\]
где:
\begin{itemize}
    \item $k_{\Omega}$ измеряет объём допустимых состояний,
    \item $k_{A}$ считает реализованные оси относительно мета-пространства,
    \item $k_{C}$ отражает устойчивость циклов,
    \item $k_{J}$ описывает когерентность потоков,
    \item $k_{T}$ измеряет структурное напряжение относительно порогов.
\end{itemize}


% ================================================================
\subsection{Иерархия уровней $K_0 \to K_{12}$}
% ================================================================

Каждый уровень представляет собой отдельную стратификацию организации:

\begin{itemize}
    \item $K_0$ — метадавление, условия возможности и препредметная логика;
    \item $K_1$ — непрерывные одномерные континуумы;
    \item $K_2$ — двумерные континуумы с внутренними потоками и кластерами;
    \item $K_3$ — химические конфигурационные пространства и многообразия реакций;
    \item $K_4$ — континуумы уровня протоклеток с мембранами и границами;
    \item $K_5$ — нейронные континуумы с циклами возбуждения и аттракторами;
    \item $K_6$ — когнитивные континуумы с предиктивными моделями;
    \item $K_7$ — социальные континуумы с групповыми устойчивыми циклами;
    \item $K_8$ — цивилизационные континуумы с потоками крупномасштабной координации;
    \item $K_9$ — теоретические континуумы (структура объяснений);
    \item $K_{10}$ — метатеоретические континуумы (структура пространств структуры);
    \item $K_{11}$ — формальные онтологии и логические континуумы;
    \item $K_{12}$ — метаонтологические континуумы, задающие пространство
                     возможностей всех нижних уровней.
\end{itemize}

Более высокие уровни не отменяют нижние; они включают их и расширяют:
\[
K_x \subseteq M_x,
\qquad
M_x \subseteq \Omega(K_{x+1}).
\]


% ================================================================
\subsection{Границы, коллапс и смерть континуумов}
% ================================================================

Каждый $K_x$ имеет границу $\partial\Omega(K_x)$, где
континуум теряет допустимость.
Континуум «умирает», когда:
\[
\Omega(K_x) = \varnothing
\quad\Rightarrow\quad
k(K_x) = 0,
\]
что соответствует:
\begin{itemize}
    \item нарушению порогов $\Theta(K_x)$,
    \item потере устойчивых циклов $C(K_x)$,
    \item потокам-деструкторам, достигающим критических значений,
    \item схлопыванию осей или несовместимости с $M_x$.
\end{itemize}

Смерть континуума отличается от эволюции: она убирает континуум из
пространства допустимых структур и не может быть обращена.


% ================================================================
\subsection{Роль \texorpdfstring{$k$}{k}-уровневой структуры в ядре модели}
% ================================================================

$K$-уровни составляют вертикальный каркас Онтологии Континуумов.
Они:
\begin{itemize}
    \item организуют зарождение от физических структур к когнитивным,
          социальным и метаонтологическим структурам;
    \item определяют рекурсию $K_x \to M_x \to K_{x+1}$;
    \item обеспечивают размерностную лестницу, необходимую для универсальной
          применимости;
    \item ограничивают эволюцию континуумов через условия допустимости;
    \item объединяют все области реальности под одной математической рамкой.
\end{itemize}

Структура $K$-уровней делает онтологию иерархической и рекурсивной:
каждый уровень зависит от мета-пространства над ним и ограничивает уровень
ниже.
